% !TEX root = ../print.tex
% Draft \S7 --- Self-decomposability and the characterization theorem
%
% Source: draft/spatial-hemigroup-scale-space.md, \S7 (lines 365-478).
%
% DRAFT NUMBERING (reported, not corrected here): the draft carries TWO statements numbered
% 7.6' -- "Lemma 7.6' (the admissible exponents form a convex cone)" and "Remark 7.6' (the
% scale evolution)" -- and states Lemma 7.6' before Proposition 7.5'. The nodes below use
% labels, so nothing here depends on the collision; the draft should renumber.
% (Checked at the draft sync, 2026-09-11: the draft had already renumbered -- the scale-evolution
% remark is "Remark 7.7'" -- so the paragraph above records a past state of the draft; the
% ordering note, Lemma 7.6' before Proposition 7.5', is still true.)
%
% lem:selfdecomposable-exponents is the analytic heart. Its (1) => (3) leg is proved
% DIRECTLY, from the dilation identity, the uniqueness of the pair, and a convex-tail
% argument in the log-displacement coordinate; that is why the classical one-dimensional
% self-decomposability theorem (prop:sd-exponents, ledger A7) is cited for orientation and
% not consumed. The causal article takes the opposite route and carries the corresponding
% classical statement in its trust base (its ledger A18), so this is a place where the
% spatial development is NARROWER than its twin.
%
% prop:strict-positivity and prop:kernel-regularity are the draft's Proposition 7.5' SPLIT:
% clauses (1) and (2) are elementary and stay [T]; clause (3), absolute continuity, and the
% unimodality claim the draft makes in the paragraph after it are classical and become one
% [A] node. The split keeps two elementary facts out of the trust base.

\chapter{Self-decomposability and the characterization theorem}
\label{ch:characterization}

By \ref{lem:additivity} the exponent of a kernel is a difference of accumulated exponents,
and by \ref{prop:canonical-gauge} the accumulated exponent is a dilate of one function. So
every kernel is read off $F$ by stretching its argument and subtracting,
\[
  g_{s,t}(\omega) = F(\tilde t\,\omega) - F(\tilde s\,\omega),
  \qquad \tilde s = \chi(s) \le \tilde t = \chi(t),
\]
and since $\chi$ is onto, $(\tilde s,\tilde t)$ ranges over all pairs $0 \le \tilde s \le
\tilde t$. \ref{thm:increments-levy} requires each of these differences to be a symmetric \Levy{}
exponent. That is a condition on $F$ alone --- not that $F \in \LEs$, which it is, but that
every stretched difference is --- and it is the condition this chapter solves. Canonical-gauge
scales are written $s \le t$, plain, when no other family is in play, so that $a$ is free for
the Gaussian coefficient.

The condition has a probabilistic name. Let $X_1$ be the random displacement whose law is the
kernel at canonical scale $1$, so that $\mathbb{E}\cos(\omega X_1) = e^{-F(\omega)}$; in the
canonical gauge the kernel at scale $t$ is the law of $tX_1$. That $F(\omega) - F(c\omega)$ is
a \Levy{} exponent for $c \in (0,1)$ says that there is a random variable $R_c$, independent of
$X_1$, with $X_1 \seq cX_1 + R_c$. That is self-decomposability
(\ref{def:self-decomposable}): covariance turns a shrinking of the scale axis into a shrinking
of the displacement, and the cascade requires what remains to be a measurement in its own
right.

% draft: Lemma 7.1'
% twin: hcs lem:selfdecomposable-exponents (ledger A18 there, for the analysis direction).
% Deliberately NOT a reuse: the analysis direction is proved here, so this article's trust
% base does not contain the corresponding classical statement.
% CHANGED (proof of record 2026-09-09): the (3) => (1) paragraph's finiteness step is
% rewritten to the machine-checked route. It said the (1 wedge u^2)-integral of k(u/t) is
% finite "because F(t.) is in LEs and lem:profile-integrability applies", which assumes what
% is to be shown -- knowing F(t.) is in LEs is knowing the dilated measure satisfies the \Levy{}
% condition. The Lean proof instead runs the SAME substitution u = tx on that integral and
% compares truncations, 1 wedge t^2x^2 <= max(1,t^2)(1 wedge x^2)
% (SpatialLine.min_one_sq_mul_le); the printed proof now does the same. The closing sentence
% also no longer appeals to prop:fourier-toolbox(3): LEs is named by its representation in
% this article, and the increment is exhibited in that form, so the direction spends nothing.
% Statement unchanged.
% CHANGED (proof of record 2026-09-09, wave 3): the "measure varpi" paragraph is rewritten to
% the machine-checked route (SpatialLine.sd_exponents_profile_measure), and lem:cin-rays joins
% the \uses list. It said "with k right-continuous, nonincreasing and k(infinity) = 0, the set
% function varpi := -dk is a positive measure with k(x) = varpi((x,infinity))", which assumes a
% right-continuous representative the previous paragraph does not produce -- a profile that is
% only nonincreasing has none this development can name. The checked route builds the measure
% FIRST, as the quantile transform lem:cin-rays(2) already attaches to any nonincreasing profile
% vanishing at infinity, and reads the right-continuous modification off its tails; the tail
% identity is then an identity at every point, not almost everywhere. The two Tonelli identities
% are replaced by layer-cake INEQUALITIES with continuous weights, u^2/(1+u^2) and
% log(1+u^2)/2 in place of (u wedge 1)^2/2 and log_+ u: the blueprint's weights have
% discontinuous densities (2t.1_{t<1} and t^{-1}1_{t>1}) and both conditions are finiteness
% assertions, so nothing is lost and two indicator functions are not carried. Statement
% unchanged. This also closes finding F6, which had priced the paragraph as downstream of an
% unwritten Stieltjes construction: no Stieltjes function is needed, and the construction it
% does need was written in chapter 8 in wave 2.
% CHANGED (proof of record 2026-09-09, wave 3): the convex-tail paragraph of "(1) implies (3)"
% is rewritten to the machine-checked route (SpatialLine.exists_antitone_density). It said that
% M(theta) = nu~((theta,infinity)) is convex because its increments over intervals of fixed
% length are nonincreasing, and took k~ = -M', the right derivative. Everything after the word
% "convex" is in Mathlib; the word itself is not. What the shift property gives is
% M(theta) + M(theta + 2delta) >= 2 M(theta + delta), MIDPOINT convexity, and Mathlib has no
% "midpoint convex plus continuous implies convex" (convexOn_of_slope_mono_adjacent wants the
% adjacent-slope inequality at arbitrary triples); supplying it is the classical dyadic
% induction plus a continuity passage, and the continuity of M is itself an argument -- an atom
% of mass alpha forces uncountably many atoms of mass alpha to its left. The checked route
% defines k~ as the supremum of the dyadic difference quotients 2^n nu~((theta, theta+2^-n]),
% which is nonincreasing BY CONSTRUCTION and is a supremum of a nondecreasing sequence, and
% recovers nu~ from it by monotone convergence and one Tonelli. No convexity, no right
% derivative, no modification, and the function produced is the nonincreasing one the statement
% asks for rather than a representative of it. The rest of the direction -- the transport
% between nu and its image under log -- is not yet machine-checked; see the skeleton's
% annotation. Statement unchanged.
% CHANGED (proof of record 2026-09-09, wave 3): the "(3) implies (2)" paragraph is rewritten to
% the machine-checked route (SpatialLine.sd_exponents_three_implies_two,
% SpatialLine.sd_exponents_symbol). Three additions, no deletion.
%   (i) Cin's own calculus, which chapter 8 did not build and the printed proof assumed: the
%     integrand (1 - cos v)/v is continuous AT THE ORIGIN too, its value there being 0 and
%     |(1-cos v)/v| <= |v|/2, so Cin is C^1 on all of R and the chain rule gives
%     d/domega Cin(omega u) = (1 - cos omega u)/omega. Without that step the differentiation is
%     not justified at all.
%   (ii) the dominating constant is named -- (4/omega_0 + 2 omega_0) on (omega_0/2, 2 omega_0)
%     -- rather than left as "C on a compact subinterval", and the integrability of 1 wedge u^2
%     is credited to the varpi paragraph above, which proves it for every measure with the
%     tail property rather than only for the constructed one.
%   (iii) "F is C^1" is separated from "F is differentiable": continuity of F' is a SECOND
%     dominated-convergence argument, and the printed proof asserted the conclusion without it.
%     The extension of eq:symbol to omega <= 0 (evenness, and the factor omega at the origin)
%     is stated, since the Lean identity is at every real frequency.
%   The closing sentence now says that the direction makes no appeal to prop:fourier-toolbox --
%   LEs is named by its representation here -- which is the same observation the (3) => (1)
%   paragraph records, and the parenthesis on the causal route gains the counterexample that
%   makes "not available" a fact rather than a caution: k(x) = (log x)^{-2} at infinity is
%   admissible and has int_1^infty k(x) dx = infty.
% Statement unchanged.
% CHANGED (proof of record 2026-09-10): the "(2) implies (3)" paragraph is rewritten to the
% machine-checked route (SpatialLine.sd_exponents_two_implies_three). The printed route
% established the representation F(omega) = (1/2)a_B omega^2 + int Cin(omega u) varpi(du)
% first and then said "the integrability conditions hold because F is finite and the pair is
% unique". Formalizing showed the order has to be the other way round: the profile data cannot
% be assembled before both integrability conditions are in hand, so the representation cannot
% be the step that supplies them. The checked route therefore (a) gets the first condition from
% the \Levy{} condition on varpi alone, by the same layer cake the varpi paragraph uses, read
% backwards; (b) gets the second from int Cin d varpi <= int_0^1 F' <= F(1), an INEQUALITY,
% which is all that is needed and which makes the improper endpoint a monotone limit rather
% than an improper integral; and (c) identifies F with the profile exponent by their
% derivatives rather than by re-deriving the representation, so eq:cin-superposition is never
% re-run in this direction. Statement unchanged.
\begin{lemma}[symmetric self-decomposable exponents]\label{lem:selfdecomposable-exponents}
\uses{prop:fourier-toolbox,lem:profile-integrability,lem:cin-rays}
\lean{SpatialLine.sd_exponents_three_implies_one, SpatialLine.sd_exponents_one_implies_three,
SpatialLine.sd_exponents_three_implies_two, SpatialLine.sd_exponents_two_implies_three,
SpatialLine.sd_exponents_profile_measure, SpatialLine.sd_exponents_symbol}\leanok
\notes{self-decomposability}
For $F \in \LEs$ with pair $(a,\nu)$ the following are equivalent.
\begin{enumerate}
\item $F(t\,\cdot) - F(s\,\cdot) \in \LEs$ for all $0 < s \le t$.
\item $F$ is continuously differentiable on $(0,\infty)$ and the symbol
$B(\omega) := \omega F'(\omega)$ lies in $\LEs$.
\item $F$ has the representation \ref{eq:sd-profile}: $\nu(dx) = k(x)\,x^{-1}dx$ with
$k : (0,\infty) \to [0,\infty)$ nonincreasing, $\int_0^1 x\,k(x)\,dx < \infty$ and
$\int_1^\infty k(x)\,x^{-1}dx < \infty$.
\end{enumerate}
In (3) the profile $k$ may be taken right-continuous with $k(\infty) = 0$, and then
$\varpi := -dk$ is a positive measure on $(0,\infty)$ with $k(x) = \varpi((x,\infty))$,
$\int (1 \wedge u^2)\,\varpi(du) < \infty$ and $\int \log_+ u\,\varpi(du) < \infty$; in (2) the
pair of $B$ is $(2a,\varpi)$.
\statusT\quad\emph{(The three conditions are not equally deep, and the asymmetry organizes
what follows. (3) implies (1) is elementary --- a change of variables in the representation,
using no property of the positive definite class --- and it is all the constructive direction
of \ref{thm:main-characterization} consumes. (1) implies (3) spends the uniqueness clause of
\ref{prop:fourier-toolbox}(3) once, and then an argument about monotone measures that is
elementary but genuine; it is the analysis direction of the main theorem. (2) is here a
consequence rather than a stepping stone, because the line offers no a priori
differentiability: what the superposition \ref{eq:cin-superposition} buys instead is that
every admissible exponent is $C^1$ off the origin with a symbol that is again a \Levy{}
exponent. \textbf{Machine-checked} (2026-09-10) at all six declarations. The whole node spends
\ref{prop:fourier-toolbox}(3)'s uniqueness clause and nothing else: once in (1) implies (3), to
read the dilation identity backwards, and once in (2) implies (3), to match the pair of $F$
with the pair the profile produces. The other four declarations, including the two elementary
directions and the construction of $\varpi$, rest on Lean core.)}
\end{lemma}

\begin{proof}
\leanok
\emph{(3) implies (1): the dilation identity.} Write
$F(c\omega) = ac^2\omega^2 + \int_0^\infty (1 - \cos c\omega x)\,k(x)\,\frac{dx}{x}$ for
$c > 0$ and substitute $u = cx$. The measure $dx/x$ is invariant under dilations, so
\[
  F(c\omega) = a c^2\,\omega^2
   + \int_0^\infty \bigl(1 - \cos\omega u\bigr)\,k(u/c)\,\frac{du}{u} :
\]
dilating the argument by $c$ dilates the profile and rescales the Gaussian coefficient, and
does nothing else. Subtracting the identity at $c = s$ from the one at $c = t$,
\begin{equation}\label{eq:dilation-difference}
  F(t\omega) - F(s\omega) = a(t^2 - s^2)\,\omega^2
   + \int_0^\infty \bigl(1 - \cos\omega u\bigr)\,\bigl(k(u/t) - k(u/s)\bigr)\,\frac{du}{u},
\end{equation}
which is of the form \ref{eq:levy-khintchine}: the Gaussian coefficient $a(t^2-s^2)$ is
nonnegative since $s \le t$, and the density is nonnegative since $u/t \le u/s$ and $k$ is
nonincreasing. For the \Levy{} condition, the density is dominated by $k(u/t)$, so it is enough to
bound the $(1 \wedge u^2)$-integral of $k(u/t)\,u^{-1}du$; and the same substitution $u = tx$
turns that integral into $\int_0^\infty (1 \wedge t^2x^2)\,k(x)\,x^{-1}dx$, at most
$\max(1,t^2)\int_0^\infty (1 \wedge x^2)\,k(x)\,x^{-1}dx$, which is finite by
\ref{lem:profile-integrability}; the comparison $1 \wedge t^2x^2 \le \max(1,t^2)(1 \wedge x^2)$
holds in each of the two regimes of the truncation separately. So the difference is itself of
the form \ref{eq:levy-khintchine}, hence in $\LEs$; the class being named by its
representation, this direction makes no appeal to \ref{prop:fourier-toolbox}.

% CHANGED (2026-09-11, post-sync follow-up, R109): the paragraph heading still said "and the
% convex tail", which the wave-3 proof-of-record rewrite (the block above) removed from the
% argument: the density is the supremum of dyadic difference quotients and no convexity is used.
% Heading only; no statement or proof changed.
\emph{(1) implies (3): the dilation identity read backwards, and the monotone density.} Fix
$c \in (0,1)$. By hypothesis $F(\omega) - F(c\omega) = a'\omega^2 + \int (1-\cos\omega
x)\,\nu'(dx)$ for some $a' \ge 0$ and some \Levy{} measure $\nu' \ge 0$. On the other hand, with
$\Dil{c}\nu$ the image of $\nu$ under $x \mapsto cx$, the first display of this proof reads
$F(c\omega) = ac^2\omega^2 + \int (1-\cos\omega x)\,(\Dil{c}\nu)(dx)$. Hence
\[
  a'\omega^2 + \int \bigl(1 - \cos\omega x\bigr)\,(\nu' + \Dil{c}\nu)(dx)
  \;=\; a\,\omega^2 + \int \bigl(1 - \cos\omega x\bigr)\,\nu(dx)
  \qquad \text{for every } \omega,
\]
an identity between two representations of the form \ref{eq:levy-khintchine} with
\emph{positive} data; so by the uniqueness clause of \ref{prop:fourier-toolbox}(3),
$a' = a(1-c^2)$ and $\nu' + \Dil{c}\nu = \nu$. Thus
\begin{equation}\label{eq:dilation-decrease}
  \Dil{c}\nu \ \le\ \nu \qquad \text{for every } c \in (0,1).
\end{equation}
Pass to the log-displacement coordinate $\theta = \log x$: let $\tilde\nu$ be the image of
$\nu$ under $\log$, a measure on $\RR$ with
$\tilde\nu((\theta,\infty)) = \nu((e^\theta,\infty)) < \infty$ for every $\theta$, finiteness
coming from $\int (1 \wedge x^2)\,\nu < \infty$. Dilation by $c$ becomes translation by
$\log c < 0$, so \ref{eq:dilation-decrease} reads $\tilde\nu(A + h) \le \tilde\nu(A)$ for every
Borel $A$ and every $h > 0$: the measure decreases under every right shift. Put
\[
  \tilde k(\theta) \;:=\; \sup_{n \ge 0}\ 2^n\,\tilde\nu\bigl((\theta,\ \theta + 2^{-n}]\bigr).
\]
Each term is nonincreasing in $\theta$, by the shift property; and the sequence is nondecreasing
in $n$, because splitting $(\theta, \theta+2\delta]$ at its midpoint and applying the shift
property to the right half gives
$\tilde\nu((\theta,\theta+2\delta]) \le 2\,\tilde\nu((\theta,\theta+\delta])$. So $\tilde k$ is
nonincreasing, and it is a supremum of a nondecreasing sequence. Fix $a < b$. By monotone
convergence and then Tonelli,
\[
  \int_a^b \tilde k \;=\; \sup_n\ 2^n\!\int_a^b \tilde\nu\bigl((\theta,\theta+\varepsilon_n]\bigr)
    d\theta
  \;=\; \sup_n\ 2^n\!\int \bigl|(a,b] \cap [x - \varepsilon_n,\, x)\bigr|\ \tilde\nu(dx),
  \qquad \varepsilon_n = 2^{-n},
\]
and the inner length is $\varepsilon_n$ for $x \in (a+\varepsilon_n, b]$, at most
$\varepsilon_n$ always, and $0$ outside $(a, b+\varepsilon_n]$. So the $n$-th term lies between
$\tilde\nu((a+\varepsilon_n, b])$ and $\tilde\nu((a, b+\varepsilon_n])$, and both converge to
$\tilde\nu((a,b])$ --- the first by continuity of a measure from below, the second by continuity
from above, which is where the finiteness of the rays is used. Hence
$\tilde\nu((a,b]) = \int_a^b \tilde k$ for every $a < b$, so $\tilde\nu(d\theta) =
\tilde k(\theta)\,d\theta$. The supremum defining $\tilde k$ is extended-valued at first, and
that identity is what makes it finite: a nonincreasing function taking the value $+\infty$ at
one point takes it at every point to the left, so some $\int_a^b \tilde k$ would be infinite,
whereas every $\tilde\nu((a,b])$ is finite. Returning to $x$, $\nu(dx) = k(x)\,x^{-1}dx$ with
$k(x) := \tilde k(\log x)$ nonincreasing; and $k(\infty) = 0$, since $\tilde k$ is nonincreasing
and nonnegative with $\int_\theta^\infty \tilde k = \tilde\nu((\theta,\infty)) < \infty$. The two
integrability
conditions are $\int (1 \wedge x^2)\,\nu(dx) < \infty$ read through
\ref{lem:profile-integrability}.

\emph{The measure $\varpi$, and the right-continuous profile.} A profile that is merely
nonincreasing has no right-continuous representative to hand, so the measure is built first and
the modification read off it. Take for $\varpi$ the measure \ref{lem:cin-rays}(2) attaches to
$k$: the image of Lebesgue measure on $(0,\infty)$ under the generalized inverse
$y \mapsto \sup\{u > 0 : k(u) > y\}$, whose tails satisfy $\varpi((x,\infty)) = k(x)$ at every
continuity point of $k$, hence at almost every $x > 0$. Every positive tail is then finite,
since a point of that identity below $x$ bounds $\varpi((x,\infty))$ by a value of $k$; so
$k^{\ast}(x) := \varpi((x,\infty))$ is a finite nonincreasing function on $(0,\infty)$, equal to
$k$ almost everywhere and therefore inducing the same \Levy{} measure, right-continuous because
$(x,\infty)$ is the increasing union of the $(x + 1/n,\infty)$, and vanishing at infinity
because $k$ does. This is the modification the statement asserts, and its tail identity holds at
\emph{every} positive point rather than almost everywhere.

For the two integrability conditions, the layer cake formula against $\varpi$ turns
$\int G(u)\,\varpi(du)$, for $G(u) = \int_0^u g(t)\,dt$ with $g \ge 0$ continuous, into
$\int_0^\infty k^{\ast}(t)\,g(t)\,dt$. Take $g(t) = 2t(1+t^2)^{-2}$, so that
$G(u) = u^2/(1+u^2) \ge \tfrac12 (1 \wedge u^2)$, and note $g(t) \le 2t$ and
$g(t) \le 2t^{-1}$; then
$\int (1 \wedge u^2)\,\varpi(du) \le 4\int_0^1 t\,k(t)\,dt + 4\int_1^\infty k(t)\,t^{-1}dt$.
Take $g(t) = t(1+t^2)^{-1}$, so that $G(u) = \tfrac12\log(1+u^2) \ge \log_+ u$, with the same
two bounds on $g$; then $\int \log_+ u\,\varpi(du)$ is bounded by the same two integrals. So
both conditions on $\varpi$ follow from $\int_0^1 x\,k(x)\,dx < \infty$ and
$\int_1^\infty k(x)\,x^{-1}dx < \infty$.

\emph{(3) implies (2): the $\Cin$ superposition.} Tonelli over $\{0 < x < u\}$ turns
\ref{eq:sd-profile} into
\begin{equation}\label{eq:cin-superposition}
  F(\omega) = a\,\omega^2 + \int_{(0,\infty)} \Cin(\omega u)\,\varpi(du),
  \qquad \Cin(z) := \int_0^{z}\frac{1 - \cos v}{v}\,dv,
\end{equation}
since $\int_0^u (1-\cos\omega x)\,x^{-1}dx = \Cin(\omega u)$; the measure $\varpi$ is the one
supplied above, and only its almost-everywhere tail identity is used. The integrand
$v \mapsto (1-\cos v)/v$ of $\Cin$ is continuous at the origin as well, its value there being
$0$ and $|(1-\cos v)/v| \le |v|/2$; so $\Cin$ is continuously differentiable on all of $\RR$
with $\Cin'(v) = (1-\cos v)/v$, and
$\partial_\omega \Cin(\omega u) = u\,\Cin'(\omega u) = (1 - \cos\omega u)/\omega$. For $\omega$
in the interval $(\omega_0/2, 2\omega_0)$ this is bounded by
$(4\omega_0^{-1} + 2\omega_0)(1 \wedge u^2)$ --- it is at most $2/\omega$ and at most
$\omega u^2/2$, which are the two regimes of the truncation --- and $1 \wedge u^2$ is
$\varpi$-integrable, as the previous paragraph shows for every such $\varpi$. So
differentiation under the integral sign applies at every $\omega > 0$, and
\begin{equation}\label{eq:symbol}
  B(\omega) = \omega F'(\omega)
   = 2a\,\omega^2 + \int_{(0,\infty)} \bigl(1 - \cos\omega u\bigr)\,\varpi(du),
\end{equation}
which is of the form \ref{eq:levy-khintchine} with pair $(2a,\varpi)$, hence in $\LEs$ --- the
class being named by its representation, so exhibiting the pair is the conclusion and no appeal
to \ref{prop:fourier-toolbox} is made. That $F$ is $C^1$ on $(0,\infty)$ and not merely
differentiable is a second dominated-convergence argument: $1 - \cos\omega v$ is bounded by
$(2 + (|\omega_0|+1)^2)(1 \wedge v^2)$ for $\omega$ within $1$ of $\omega_0$, so
$\omega \mapsto \int (1 - \cos\omega v)\,\varpi(dv)$ is continuous on $\RR$ and $F' $ is
continuous where $\omega \ne 0$. Both sides of \ref{eq:symbol} are stated at every real
$\omega$: the identity extends to $\omega < 0$ because $F$ is even, hence $F'$ odd, and holds
at $\omega = 0$ because the factor $\omega$ annihilates whatever value the derivative is
assigned there. (The half-line route, differentiating under the integral in the profile form and
applying Tonelli, is not available here: $\omega\sin(\omega x)$ is not of one sign, and
$\int u\,\varpi(du)$ may diverge, as it does for the stable profiles of index below $1$. The
$\Cin$ form has a positive integrand and needs no such integral. Nor is the profile form saved
by the two admissibility conditions: $k(x) = (\log x)^{-2}$ at infinity satisfies
$\int_1^\infty k(x)x^{-1}dx < \infty$ while $\int_1^\infty k(x)\,dx$ diverges, and that
integral is what dominating $|\sin(\omega x)|k(x)$ would need.)

\emph{(2) implies (3).} Let $B = \omega F'$ have pair $(a_B,\varpi)$ and put
$k(x) := \varpi((x,\infty))$, which is nonnegative and nonincreasing on $(0,\infty)$ and finite
at every positive $x$, since $\int (1 \wedge u^2)\,\varpi(du) < \infty$ bounds
$\varpi((x,\infty))$ by $(1 \wedge x^2)^{-1}$ times that integral. The claim is that $F$ is
\ref{eq:sd-profile} at the data $(\tfrac12 a_B, k)$, and the two integrability conditions on
$k$ come from two different places --- which is the order the proof has to be read in, because
the data is not admissible until both are in hand.

The first is the \Levy{} condition on $\varpi$ alone. The layer cake formula of the previous
paragraph, read from right to left at $g(t) = 2t(1+t^2)^{-2}$, turns
$\int_0^\infty k(t)\,g(t)\,dt$ into $\int u^2(1+u^2)^{-1}\,\varpi(du)$, at most
$\int (1 \wedge u^2)\,\varpi(du)$; and $g(t) \ge t/2$ on $(0,1)$, so
$\int_0^1 t\,k(t)\,dt \le 2\int (1 \wedge u^2)\,\varpi(du) < \infty$.

The second is not implied by the \Levy{} condition: $\varpi(du) = du/(u\log^2 u)$ on $(2,\infty)$
has finite truncated integral and divergent logarithmic one. What supplies it is the finiteness
of $F$. Since $B \ge 0$, $F' \ge 0$ on $(0,\infty)$, so the fundamental theorem of calculus on
$[\varepsilon,1]$ gives $\int_\varepsilon^1 F'(v)\,dv = F(1) - F(\varepsilon) \le F(1)$ for
every $\varepsilon \in (0,1)$, and letting $\varepsilon \downarrow 0$ along a sequence ---
monotonically, the integrand being nonnegative --- $\int_0^1 F'(v)\,dv \le F(1) < \infty$.
Now $F'(v) \ge v^{-1}\int (1 - \cos vu)\,\varpi(du)$, and Tonelli over
$(0,1) \times (0,\infty)$ with a positive integrand, together with the substitution
$\int_0^1 (1 - \cos uv)\,v^{-1}dv = \Cin(u)$, turns the integral of that lower bound into
$\int \Cin(u)\,\varpi(du)$. So $\int \Cin(u)\,\varpi(du) \le F(1) < \infty$. Since
$\Cin \ge 0$ everywhere and $\Cin(u) = \log u + O(1)$ at infinity (\ref{lem:cin-rays}(1)),
$\log_+ u \le \Cin(u) + C\,(1 \wedge u^2)$ for a fixed $C$, so
$\int \log_+ u\,\varpi(du) < \infty$; the second layer cake, at $g(t) = t(1+t^2)^{-1}$, which
is at least $(2t)^{-1}$ beyond $1$, converts that into $\int_1^\infty k(t)\,t^{-1}dt < \infty$.

So $(\tfrac12 a_B, k)$ is admissible data, and defines an exponent $\tilde F$ of the form
\ref{eq:sd-profile}. It is $F$, and the identification does not re-derive the representation:
by \ref{eq:symbol} applied to $\tilde F$, whose profile has $\varpi$ for its tail measure,
$\omega\tilde F'(\omega) = a_B\omega^2 + \int (1 - \cos\omega u)\,\varpi(du) = B(\omega)
 = \omega F'(\omega)$ at every real $\omega$, so $\tilde F$ and $F$ have the same derivative at
every $\omega > 0$; both are continuous and vanish at the origin, so they agree on $[0,\infty)$
--- a function continuous on $[0,\infty)$ with vanishing derivative on $(0,\infty)$ is constant
on every $[\varepsilon, \omega]$, and its value at the origin is the limit of those constants
--- and both are even, so they agree on $\RR$. The uniqueness clause of
\ref{prop:fourier-toolbox}(3), applied to the two representations of that one function, then
identifies the pairs: $a = \tfrac12 a_B$ and $\nu(dx) = k(x)\,x^{-1}dx$, which is (3).
\end{proof}

% CHANGED (article mirror 2026-09-14, R157): two wordings and one anchor. The increments are "not
% NECESSARILY stationary" -- the Cauchy member has stationary increments in the canonical gauge,
% so the old phrase excluded a member of the class (the external review's finding on this remark);
% and a is not a variance RATE, the continuous part having variance 2a tilde-t^2, which is
% quadratic in the canonical gauge. The causal sentence is dropped, the comparison belonging to
% Chapter~\ref{ch:bridge}.
\begin{remark}[the probabilistic reading]\label{rem:sato-process}
Together with the similarity structure, \ref{eq:sd-profile} says that the kernels are the
marginal laws of a \emph{symmetric self-similar additive process}: a process
$(X_{\tilde t})_{\tilde t \ge 0}$ on $\RR$ with independent, symmetric, not necessarily
stationary increments satisfying $X_{\lambda\tilde t} \seq \lambda X_{\tilde t}$, the Sato
process of exponent $1$ attached to the self-decomposable law of $X_1$
(\sscite{sato1999levy}, Thm.~16.1, p.~99). The paths oscillate, and the Gaussian coefficient
$a$ sets the variance $2a\tilde t^{\,2}$ of their continuous part. Nothing in the proofs uses
this correspondence, which is stated as a reading and grounds no statement.
% CHECK (draft, checks for \S7, item 6) --- RESOLVED (2026-09-14): the anchor for the
% correspondence between self-decomposable laws and self-similar additive processes is
% Sato \S16, Thm. 16.1, p. 99 (class L is exactly the class of marginals of self-similar
% additive processes), read from the held copy by the librarian on 2026-09-14 and now printed in
% the remark. Nothing in the article's proofs uses it, so it still grounds no node and no ledger
% entry.
\end{remark}

% CHANGED (skeleton 2026-09-08): the kernel sentence read "the scale-space kernels are
% phi_t = FT^{-1}[e^{-F(t .)}], symmetric probability DENSITIES for t > 0". Absolute continuity
% of mu_{0,t} is prop:kernel-regularity, an [A] node grounded in A8, and this theorem neither
% does nor can \uses it -- the reverse edge exists and adding this one makes the graph cyclic.
% Stating it here would put A8 in the HEADLINE theorem's #print axioms and contradict its own
% [T] grade. Narrowed to "the symmetric probability measure with transform e^{-F(t .)}", which
% is what the construction proves; the density and its unimodality stay in
% prop:kernel-regularity. Found by writing the Lean statement (Skeleton.main_construction).
%
% CHANGED (review 2026-09-09): the status annotation now says so in words -- that the kernels
% have densities, that this is prop:kernel-regularity's clause and not this theorem's, and why
% the edge cannot be added here. Answering Q6: the narrowing above is invisible to a reader of
% the headline alone unless the annotation states it, and the annotation is the right place
% because it is not part of what the theorem asserts.
%
% draft: Theorem 7.3'
% twin: hcs thm:main-characterization
% (Hemigroup.SelfDecomposableExponent.main_characterization), which the causal development
% splits into thm:main-construction, thm:main-analysis and prop:main-uniqueness. The Lean
% design here is expected to make the same split; the node is kept whole because the draft
% states it whole and because the split is a formalisation decision, not a mathematical one.
% CHANGED (proof of record 2026-09-09): three steps of the sufficiency half are rewritten to
% the machine-checked route (SpatialLine.main_construction).
%   (A7): the epsilon/3 argument on a compact carrying mass 1-eps, first for compactly
%     supported continuous f and then by density, is replaced by the estimate bounding the
%     L1 distance from mu*f to f by the mu-average of the modulus of continuity of
%     translation -- a BOUNDED CONTINUOUS function of the shift vanishing at 0, because the
%     translation group acts continuously on L1. \Levy{} continuity then finishes with no
%     compact and no density step. The two-endpoint structure a hemigroup forces is made
%     explicit, since it is what the proof actually does.
%   (ND): the appeal to prop:strict-positivity(2) is replaced by the weaker fact the clause
%     needs -- that SOME frequency is moved -- which follows from lem:dilation-invariance
%     alone. The prop:strict-positivity edge in the \uses list is kept, over-approximating
%     being harmless, but the sufficiency half no longer depends on that node.
%   The gauge: 'an increasing bijection of [0,infty) is continuous' is given its reason.
% Statement unchanged.
%
% NOTE (proof 2026-09-09, wave 3): the node is CLOSED. All three clauses are proved, in
% SpatialLine/MainAnalysis.lean and SpatialLine/MainConstruction.lean, and the bundle
% SpatialLine.main_characterization names all three in one declaration, on the pattern of
% matern_theorem. Its footprint is Lean core plus A1's forward implication and A3's converse
% and uniqueness clauses -- no A8, no A9, no A11.
%   The UNIQUENESS clause has R1's four gauge hypotheses and the printed route, and prints Lean
% core, which took one new library file: lem:dilation-invariance wants continuity of F at the
% origin, and the only route to it in the library ran through prop:fourier-toolbox(3)'s
% converse, spending ledger A3 for a fact dominated convergence gives outright
% (SpatialLine/ExponentContinuity.lean).
%   The ANALYSIS direction was assembled in one worktree and its single remaining hypothesis
% proved in another: SpatialLine.main_analysis_of_profileForm is the statement with
% lem:selfdecomposable-exponents (1) => (3) as its only hypothesis, and the merge applied it to
% SpatialLine.sd_exponents_one_implies_three. The assembly's prediction -- 'it waits on that
% declaration and on nothing else' -- held exactly. The one step of the assembly worth naming
% is that the gauge's SURJECTIVITY is what makes condition (1) available at every pair of
% canonical scales; it is the only place in the chapter that clause is consumed.
\begin{theorem}[main theorem]\label{thm:main-characterization}
\uses{def:cascade-family,lem:convolution-representation,lem:nonvanishing,lem:additivity,thm:increments-levy,lem:covariance-fourier,lem:action-rigidity,prop:canonical-gauge,lem:selfdecomposable-exponents,prop:fourier-toolbox,prop:levy-continuity,prop:strict-positivity,lem:dilation-invariance,prop:fourier-uniqueness}
\lean{SpatialLine.main_characterization, SpatialLine.main_construction,
SpatialLine.main_analysis_exists, SpatialLine.main_analysis,
SpatialLine.main_uniqueness}\leanok
\notes{hemigroup-cascade}
A family $(\Phis{s}{t})_{0 \le s \le t}$ satisfies (A1)--(A8) and (ND) \emph{if and only if}
there are an increasing bijection $\chi$ of $[0,\infty)$, the gauge, and a function $F$ of the
form \ref{eq:sd-profile} with $F \not\equiv 0$, such that
\[
  \Phis{s}{t} f = \mu_{s,t} * f, \qquad
  \hat\mu_{s,t}(\omega) = \exp\bigl[-\bigl(F(\chi(t)\,\omega) - F(\chi(s)\,\omega)\bigr)\bigr].
\]
In the canonical gauge the kernel at scale $t$ is the symmetric probability measure $\phi_t$
with $\hat\phi_t = e^{-F(t\,\cdot)}$, and the pair $(\chi,F)$ is unique up to the normalization
$\chi(1) = 1$.
\statusT\quad\emph{(The headline. Read against the axioms: the apertures are the marginal laws
of one symmetric self-similar additive process, the kernel at canonical scale $t$ being the law
of $tX_1$, so the whole family is generated by one law and the gauge. What the axioms leave
free is exactly the pair $(a,k)$ of \ref{eq:sd-profile}: a jitter variance and a nonincreasing
displacement profile. The Gaussian scale space is $(a,0)$; everything else has a catalogue. The
formalisation splits the statement three ways --- construction, analysis, uniqueness --- on the
causal development's model, and the three are assembled in one declaration,
\texttt{main\_characterization}, so that the \texttt{lean} tag names all three clauses. The
necessity direction appears there in two forms, and both are named. One concludes the
representation --- from the axioms alone there \emph{are} kernels $\mu_{s,t}$ with
$\Phis{s}{t}f = \mu_{s,t} * f$, together with the gauge and the profile --- which is the
sentence this theorem states; the other takes a representing family as a hypothesis, which is
the shape every lemma of Chapters 5 to 7 uses and which says more about a \emph{given} family,
since it quantifies over all of them. Nothing is lost by concluding the representation, because
\ref{lem:convolution-representation} determines the kernel at each admissible pair uniquely.
Symmetry of the kernels is a conjunct of \emph{both} directions, not of the sufficiency
direction alone: $\hat\mu_{s,t}$ is the cosine transform, which is the measure's characteristic
function only where the measure is symmetric, so a necessity clause that omitted the rider would
say less about $\mu_{s,t}$ than the display appears to. It is (A3) read through the uniqueness
half of \ref{lem:convolution-representation}, and costs one application of it. The
normalization is carried across the join: the necessity clause concludes $\chi(1) = 1$, which
is what the uniqueness clause quantifies over, so the two meet. That was worth stating rather
than leaving to the reader, because the rest of the necessity conclusion is invariant under
$(\chi,F) \mapsto (c\,\chi,\ F(c^{-1}\,\cdot))$ for every $c > 0$ and singles out no member of
that orbit. Its
\texttt{\#print axioms} is Lean core together with two ledger entries and no more:
\ledger{A1}, Bochner's symmetric leg in the forward direction, and \ledger{A3}, the symmetric
\Levy--Khintchine representation at its converse and its uniqueness clauses. The kernels
have densities as well, unimodal about the origin, but that is \ref{prop:kernel-regularity} and
not this theorem: the statement here is about measures, deliberately, absolute continuity
resting on an interface this theorem cannot cite without making the dependency graph cyclic.
A reader who wants the density should read the two nodes together.)}
\end{theorem}

\begin{proof}
\leanok
\emph{Necessity.} \ref{lem:convolution-representation}, \ref{lem:nonvanishing},
\ref{lem:additivity}, \ref{thm:increments-levy}, \ref{lem:covariance-fourier},
\ref{lem:action-rigidity} and \ref{prop:canonical-gauge} together give the representation with
$F = G(1,\cdot) \in \LEs$ and $F \not\equiv 0$. Every increment
$F(\chi(t)\,\cdot) - F(\chi(s)\,\cdot) = g_{s,t}$ lies in $\LEs$ by \ref{thm:increments-levy},
and $\chi$ is onto $(0,\infty)$ by \ref{prop:canonical-gauge}, so condition (1) of
\ref{lem:selfdecomposable-exponents} holds for every pair $0 < \tilde s \le \tilde t$ of
canonical scales, and that lemma yields \ref{eq:sd-profile}. The gauge so obtained is already
\emph{normalized}: it is the inverse of the orbit coordinate $c(\lambda) = S_\lambda 1$, which
fixes $1$ by \ref{prop:canonical-gauge}, so $\chi(1) = 1$. What this direction returns is
therefore a normalized representation --- one of those the uniqueness clause compares --- and
not merely a representative of its orbit under
$(\chi,F) \mapsto (c\,\chi,\ F(c^{-1}\,\cdot))$.

\emph{Sufficiency.} Let $F$ be of the form \ref{eq:sd-profile} with $F \not\equiv 0$, and work
in the canonical gauge. For $0 \le s \le t$ the function
$g_{s,t}(\omega) := F(t\omega) - F(s\omega)$ lies in $\LEs$ by
\ref{lem:selfdecomposable-exponents}, (3) implies (1) --- for $s = 0$ the increment is
$F(t\,\cdot)$ itself, which is in $\LEs$ by hypothesis and the dilation identity --- so by the
converse half of \ref{prop:fourier-toolbox}(3) and by \ref{prop:fourier-toolbox}(1),
$e^{-g_{s,t}}$ is the transform of a unique probability measure $\mu_{s,t}$ on $\RR$, symmetric
because the transform is real and even. Put $\Phis{s}{t}f := \mu_{s,t} * f$. Then (A1)--(A5)
hold by the converse direction of \ref{lem:convolution-representation}; (A6) because the
exponents add and transforms determine measures (\ref{prop:fourier-uniqueness}); (A8) with
$S_\lambda t = \lambda t$, since $\hat\mu_{\lambda s,\lambda t}(\omega) =
\hat\mu_{s,t}(\lambda\omega)$; and (ND) because $F(t\,\cdot) - F(s\,\cdot) \not\equiv 0$ for
$s < t$: at $s = 0$ that is $F \not\equiv 0$, and for $s > 0$, if the increment vanished
identically then $F(\omega) = F(c\omega)$ for every $\omega$ with $c = s/t \in (0,1)$, whence
$F \equiv 0$ by \ref{lem:dilation-invariance}. Only that \emph{some} frequency is moved is
needed here, which is less than \ref{prop:strict-positivity}(2) asserts.

For (A7), a collapsing-increment estimate does the work. For a probability measure $\mu$ and
$f \in \Lone$ the $\Lone$ distance from $\mu * f$ to $f$ is at most $\int \Theta_f(y)\,\mu(dy)$,
where $\Theta_f(y)$ is the $\Lone$ distance from $\Tr{y}f$ to $f$; and $\Theta_f$ is bounded,
continuous, and vanishes at the origin, because the translation group acts continuously on
$\Lone$. So if $(s_n,t_n) \to (c,c)$ along admissible pairs then $\hat\mu_{s_n,t_n} \to 1$
pointwise, by continuity of $F$ and of $\chi$; hence $\mu_{s_n,t_n} \to \delta_0$ weakly by
\ref{prop:levy-continuity}, hence $\int \Theta_f(y)\,\mu_{s_n,t_n}(dy) \to 0$: a collapsing increment
tends to the identity in the strong operator topology. At a pair $s < t$ the two endpoints are
then moved separately --- there is no $t - s$ to move, this being a hemigroup: the left endpoint
through an increment applied \emph{first}, where the contraction property absorbs it, and the
right endpoint through one applied \emph{last}, where the same estimate is read at the image.
No compact set, no density argument and no $\varepsilon/3$ are needed.

Finally, for an arbitrary gauge $\chi$ set $\Phi'_{s,t} := \Phi_{\chi(s),\chi(t)}$: (A1)--(A6)
and (ND) hold verbatim, (A7) because an increasing bijection of $[0,\infty)$ is continuous ---
it is an order isomorphism of $[0,\infty)$ onto itself, and an order isomorphism between linear
orders carrying the order topology is a homeomorphism --- and (A8) with
$S'_\lambda = \chi^{-1}(\lambda\,\chi(\cdot))$.

\emph{Uniqueness.} The exponents $g_{0,t}$ are determined by the family: the kernels by
\ref{lem:convolution-representation}, their transforms by \ref{prop:fourier-uniqueness}, and
$-\log$ is single-valued on them by \ref{lem:nonvanishing}. The normalization $\chi(1) = 1$
recovers $F = g_{0,1}$; without it the pair is determined up to
$(\chi,F) \mapsto (c\,\chi,\ F(c^{-1}\,\cdot))$ with $c > 0$, which is the freedom the
normalization removes. For each $t$, $\chi(t)$ is then the unique $c$ with
$g_{0,t} = F(c\,\cdot)$: such a $c$ exists by the representation, and if
$F(c\omega) = F(c'\omega)$ for every $\omega$ with $c/c' \ne 1$, then $F \equiv 0$ by
\ref{lem:dilation-invariance}, a contradiction.
\end{proof}

% draft: Corollary 7.4'
% twin: hcs cor:semigroup-case (Hemigroup.CascadeCore.semigroup_case)
% CHANGED (proof of record 2026-09-09, wave 3): the whole proof is rewritten to the
% machine-checked route (SpatialLine.semigroup_case, SpatialLine.semigroup_case_profile).
% Three differences, and each of them changes what the proof cites.
%   (i) The printed proof says "homogeneity gives g_{s,t} = (t-s)g" in one clause. That is a
%     Cauchy equation in the SCALE variable and it is the longest step of the checked proof:
%     one-parametricity gives equality of OPERATORS, the uniqueness half of
%     lem:convolution-representation turns it into equality of KERNELS, lem:additivity makes
%     the accumulated exponent additive in the scale, and an additive function continuous on
%     [0,infty) is linear there only after an odd extension to R. Three labels are added to
%     \uses for it: lem:convolution-representation, lem:additivity, thm:increments-levy.
%   (ii) lem:action-rigidity(3) is NOT consumed. The normalisation makes S_lambda(1) = g(lambda),
%     so the multiplicative function is g itself and there is no separate homomorphism c whose
%     continuity has to be established. What IS consumed instead is lem:no-lattice, for strict
%     positivity of g off the origin, which the logarithm needs; and lem:nonvanishing, for
%     continuity of g. Both are added to \uses; lem:action-rigidity is kept, over-approximating
%     being harmless.
%   (iii) The alpha <= 2 step is made explicit rather than asymptotic, and the profile's two
%     integrability conditions are the SAME two comparisons that make C_alpha converge, so the
%     appeal to lem:profile-integrability is dropped from the text (the label stays in \uses).
% Statement unchanged.
\begin{corollary}[recovering the semigroup case]\label{cor:semigroup-case}
\uses{thm:main-characterization,lem:convolution-representation,lem:nonvanishing,lem:additivity,thm:increments-levy,lem:covariance-fourier,lem:action-rigidity,lem:no-lattice,lem:quadratic-growth,lem:profile-integrability}
\lean{SpatialLine.semigroup_case, SpatialLine.semigroup_case_profile,
SpatialLine.semigroup_case_gaussian}\leanok
\notes{covariance-strata}
Assume in addition that the family is one-parameter, that is, $\Phis{s}{t}$ depends only on
$t - s$. Then, after the normalization $g_{0,1}(1) = 1$,
\[
  F(\omega) = |\omega|^\alpha \quad \text{for some } 0 < \alpha \le 2 :
\]
the kernels are exactly the symmetric stable densities, which are the members of the extended
semigroup class of \sscite{pauwels1995extended} with positive kernels. Here $\alpha = 2$ is the
Gaussian, and $0 < \alpha < 2$ has Gaussian coefficient $0$ and profile
$k(x) = C_\alpha^{-1}x^{-\alpha}$ with
$C_\alpha := \int_0^\infty (1 - \cos u)\,u^{-1-\alpha}du \in (0,\infty)$.
\statusT\quad\emph{(The semigroup case, recovered from inside the hemigroup theory. The bound
$\alpha \le 2$ is \ref{lem:quadratic-growth} and not a cited growth estimate, which is the one
reroute in this chapter.)}
\end{corollary}

\begin{proof}
\leanok
\emph{The exponent is linear in the scale.} One-parametricity makes $\Phis{0}{t}$ and
$\Phis{s}{s+t}$ the same operator, so the uniqueness half of
\ref{lem:convolution-representation} makes $\mu_{0,t}$ and $\mu_{s,s+t}$ the same measure, and
\ref{lem:additivity} then reads $G(s+t,\omega) = G(s,\omega) + G(t,\omega)$ for all
$s,t \ge 0$: the accumulated exponent is additive in the scale. It is continuous there as well
(\ref{lem:additivity}), and an additive function continuous on $[0,\infty)$ is linear on it ---
extend it oddly to $\RR$, where it stays additive and continuous, and a continuous additive map
between real vector spaces is $\RR$-linear. So $g_{0,t} = t\,g$ with $g := g_{0,1}$.

\emph{The exponent is multiplicative in the frequency.} \ref{eq:similarity} at $t = 1$, read
through the previous paragraph on the left-hand side, is
$S_\lambda(1)\,g(\omega) = g(\lambda\omega)$ for every $\lambda > 0$ and every $\omega$. At
$\omega = 1$ the normalization $g(1) = 1$ gives $S_\lambda(1) = g(\lambda)$, so
$g(\lambda\omega) = g(\lambda)\,g(\omega)$: the multiplicative function is $g$ itself, and no
separate homomorphism has to be identified.

\emph{Cauchy.} $g$ is continuous (\ref{lem:nonvanishing}) and strictly positive off the origin
(\ref{lem:no-lattice}), so $u \mapsto \log g(e^u)$ is additive and continuous on $\RR$, hence
linear by the same theorem as above, and $g(\omega) = \omega^\alpha$ for $\omega > 0$ with
$\alpha := \log g(e)$. Evenness of the exponent extends the formula to $\omega < 0$, and
$g(0) = 0$ covers the origin once $\alpha \ne 0$ is known.

\emph{The index.} $\alpha > 0$ because $g$ is continuous at $0$ with $g(0) = 0$, while
$\omega^\alpha \ge 1$ on $(0,1]$ when $\alpha \le 0$. For the upper bound,
\ref{lem:quadratic-growth} applied to the \Levy{} pair of $g$, which \ref{thm:increments-levy}
supplies, gives a constant $C$ with $g(\omega) \le C(1+\omega^2)$; the normalization $g(1) = 1$
forces $C \ge 1/2$, and if $\alpha > 2$ then at $\omega := (2C+1)^{1/(\alpha-2)} \ge 1$ the left
side is $(2C+1)\,\omega^2 > C(1+\omega^2)$, which the bound forbids. So $\alpha \le 2$.

\emph{The profile.} When $0 < \alpha < 2$ the integral $C_\alpha$ converges, by the two regimes
of the truncation: $1 - \cos u \le u^2/2$ makes its integrand at most $u^{1-\alpha}/2$ near the
origin, integrable exactly for $\alpha < 2$, and $1 - \cos u \le 2$ makes it at most
$2u^{-1-\alpha}$ beyond $1$, integrable exactly for $\alpha > 0$. Those two comparisons
\emph{are} the two integrability conditions of \ref{eq:sd-profile} for
$k(x) = C_\alpha^{-1}x^{-\alpha}$, which is nonincreasing, with $a = 0$, so no separate appeal
is needed; and $C_\alpha > 0$ because its integrand is nonnegative and does not vanish on
$(0,1)$. The substitution $u = |\omega|x$, performed on the $[0,\infty]$-valued integral so that
integrability does not have to be revisited at each frequency, gives
$\int_0^\infty (1-\cos\omega x)\,x^{-1-\alpha}dx = C_\alpha|\omega|^\alpha$, hence the exponent
$|\omega|^\alpha$. When $\alpha = 2$, $k \equiv 0$ and $a = 1$.
\end{proof}

% CHANGED (article mirror 2026-09-14, R157): the Gaussian coefficient adds a displacement of
% variance 2at^2 to the kernel AT SCALE t and of variance 2a(t^2 - s^2) to an increment, not
% "variance 2at^2 to every kernel", which was the external review's finding on this remark; the
% word "jitter" is out; and the opening sentence, which set the scene on the half-line, is dropped
% with the rest of the causal framing (the comparison lives in Chapter~\ref{ch:bridge}).
\begin{remark}[the boundary, the Gaussian ray, and what is not there]\label{rem:boundary}
The boundary $\alpha = 2$ of the semigroup slice is the Gaussian, nothing is degenerate, and
the trace of the boundary in the general class is the Gaussian coefficient $a$, which adds an
independent Gaussian displacement of variance $2at^2$ to the kernel at scale $t$ and of
variance $2a(t^2 - s^2)$ to every increment. The members of \sscite{pauwels1995extended} beyond
$\alpha = 2$ have positive transfer functions and sign-changing kernels. They are excluded by
(A4) through \ref{thm:main-characterization}, and by nothing in
Chapters~\ref{ch:representation} to \ref{ch:covariance}. So positivity, rather than
recursivity, narrows the class. One axiom does not transfer: the additional axiom that removed
the drift in the causal theory has a counterpart here,
$\lim_{\omega \to \infty} F(\omega)/\omega^2 = 0$, which says $a = 0$. Unlike its causal twin
it would exclude the Gaussian, and this article does not impose it.
\end{remark}

% draft: Lemma 7.6' (the FIRST statement so numbered in the draft; see the header note)
% twin: hcs lem:admissible-cone (Hemigroup.SelfDecomposableExponent.admissible_cone)
% CHANGED (proof of record 2026-09-09): the proof now says why the two integrability
% conditions and the exponent are additive with no side condition (nonnegative integrands,
% so the identity is one in [0,infty]) and where the finiteness that turns it into an
% identity between REAL numbers comes from (lem:profile-integrability into
% lem:quadratic-growth). The old one sentence asserted additivity of the integral without
% either, which is the only step of the machine-checked proof that is not arithmetic;
% lem:profile-integrability and lem:quadratic-growth are added to \uses accordingly.
% Statement unchanged.
\begin{lemma}[the admissible exponents form a convex cone]\label{lem:admissible-cone}
\uses{thm:main-characterization,lem:selfdecomposable-exponents,lem:profile-integrability,lem:quadratic-growth}
\lean{SpatialLine.admissible_cone}\leanok
If $F_1, F_2$ are admissible in the sense of \ref{thm:main-characterization}, of the form
\ref{eq:sd-profile} with Gaussian coefficients $a_1, a_2 \ge 0$ and nonincreasing profiles
$k_1, k_2$, then so are $F_1 + F_2$, with data $(a_1+a_2, k_1+k_2)$, and $c\,F_1$ for $c \ge 0$,
with data $(c\,a_1, c\,k_1)$. So the admissible exponents form a convex cone, and the
correspondence with the pairs $(a,k)$ is linear.
\statusT
\end{lemma}

\begin{proof}
\leanok
Each defining condition --- nonnegativity of $a$, nonnegativity and monotonicity of $k$, and
the two integrability conditions --- is stable under sums and under multiplication by $c \ge 0$,
the last two because the integrands $x k(x)$ and $k(x)/x$ are nonnegative on $(0,\infty)$, so that
the integral of a sum is the sum of the integrals with no side condition. The same additivity,
applied to \ref{eq:sd-profile} at a fixed $\omega$, gives $F_1 + F_2$ and $c\,F_1$ as the exponents of
the summed and scaled data, as identities in $[0,\infty]$; they are identities between real numbers
because \ref{lem:profile-integrability} makes each $F_i$ the exponent of a symmetric \Levy{} pair and
\ref{lem:quadratic-growth} makes that finite at every frequency.
\end{proof}

% draft: Proposition 7.5' clauses (1) and (2). Clause (3) is split off into
% prop:kernel-regularity below; see the header note.
% NEW: no causal counterpart. On the half-line the two facts below were available BEFORE the
% gauge argument, from the vanishing lemma and the monotonicity of Bernstein functions; here
% they are consequences of the classification.
% CHANGED (proof of record 2026-09-09): clause (1)'s proof is rewritten to the
% machine-checked route (SpatialLine.strict_positivity_monotone). The substitution
% u = omega x is correct but is a second change of variables on top of the one
% lem:selfdecomposable-exponents (3) => (1) already performs; reading that node at the
% single frequency omega = 1 gives the clause in one line, and clause (1) is thereby a
% corollary of a node already proved rather than a computation of its own. Statement
% unchanged.
% CHANGED (proof of record 2026-09-09, wave 3): clause (2)'s proof is rewritten to the
% machine-checked route (SpatialLine.strict_positivity_strict,
% SpatialLine.strict_positivity_consequences). Two differences.
%   (i) The rigidity step no longer covers (0,infty) by the intervals [x, kappa x]. The tail
%     integral J(eps) = int_eps^infty k(u) du/u is finite off the origin and invariant under
%     the dilation, so it has no mass on (eps, kappa eps]; antitonicity bounds the integrand
%     below there by its value at the RIGHT endpoint, which must therefore vanish; and
%     eps = x/kappa is an arbitrary point of the half-line in disguise. No covering, no
%     induction, and no almost-everywhere-to-pointwise step.
%   (ii) The lattice sentence does not consume lem:lattice-zero. A law carried by p Z has
%     cosine transform 1 at 2 pi / p, so its exponent vanishes at a nonzero frequency, which
%     clause (2) has just forbidden. The label is kept in \uses -- the annotation still refers
%     to it, and over-approximating is harmless -- but the proof no longer routes through it.
% Statement unchanged.
\begin{proposition}[strict positivity and monotonicity of the admissible exponents]
\label{prop:strict-positivity}
\uses{lem:selfdecomposable-exponents,lem:lattice-zero}
\lean{SpatialLine.strict_positivity_monotone, SpatialLine.strict_positivity_strict,
SpatialLine.strict_positivity_consequences}\leanok
Let $F$ be of the form \ref{eq:sd-profile} with $F \not\equiv 0$.
\begin{enumerate}
\item $F$ is nondecreasing on $[0,\infty)$.
\item $F(t\omega) - F(s\omega) > 0$ for all $0 \le s < t$ and every $\omega \ne 0$.
Consequently, in the canonical gauge, no increment $\mu_{s,t}$ with $s < t$ is a lattice law,
and $G(\cdot,\omega)$ is strictly increasing for every $\omega \ne 0$.
\end{enumerate}
\statusT\quad\emph{(Clause (2) is what \ref{cor:monotonicity} could not give before covariance
and \ref{lem:no-lattice} gave only for the kernels from the origin: after the classification
the exceptional set is empty. So the causal theory's strict monotonicity, and the monotonicity
in the transform variable that its function class supplied for free, both hold here --- for
the exponents the axioms select, not for the class $\NDs$, in which
\ref{lem:lattice-zero} shows they are false.)}
\end{proposition}

\begin{proof}
\leanok
(1) For $0 \le \omega_1 \le \omega_2$ the difference $F(\omega_2) - F(\omega_1)$ is the value at
frequency $1$ of the increment $F(\omega_2\,\cdot) - F(\omega_1\,\cdot)$, which lies in $\LEs$ by
\ref{lem:selfdecomposable-exponents}, (3) implies (1) --- at $\omega_1 = 0$ read as the dilate
$F(\omega_2\,\cdot)$ itself --- and a member of $\LEs$ is nonnegative. So monotonicity is the
dilation identity evaluated at one frequency, and no second change of variables is needed.

(2) By the same formula, for $\omega > 0$ --- which suffices, $F$ being even ---
\[
  F(t\omega) - F(s\omega)
  = a(t^2-s^2)\omega^2
   + \int_0^\infty (1-\cos u)\,\bigl[k(u/(t\omega)) - k(u/(s\omega))\bigr]\,\frac{du}{u},
\]
with a nonnegative integrand, reading $k(u/(s\omega)) = k(\infty) = 0$ when $s = 0$. Suppose it
vanishes. Then $a = 0$, because $t^2 > s^2$ and $\omega \ne 0$; and the integral vanishes, so
its integrand vanishes for almost every $u > 0$, and since $1 - \cos u$ is positive off the
countable set $2\pi\ZZ$, which is null, $k(u/(t\omega)) = k(u/(s\omega))$ for almost every
$u > 0$. For $s = 0$ this says $k = 0$ almost everywhere, and $F \equiv 0$.

For $s > 0$ it says $k(\varkappa x) = k(x)$ for almost every $x > 0$, with
$\varkappa := t/s > 1$, and the tail integral finishes it. For every $\varepsilon > 0$ the
integral $J(\varepsilon) := \int_\varepsilon^\infty k(u)\,du/u$ is finite: bound $k(u)/u$ by
$k(\varepsilon)/\varepsilon$ on $(\varepsilon,1]$, using that $k$ is nonincreasing, and use
$\int_1^\infty k(u)\,u^{-1}du < \infty$ beyond $1$. The change of variables $u = \varkappa x$
carries $J(\varkappa\varepsilon)$ to $\int_\varepsilon^\infty k(\varkappa x)\,dx/x$, which is
$J(\varepsilon)$ by the almost-everywhere relation --- the measure $dx/x$ being invariant under
the multiplicative group. Splitting $J(\varepsilon)$ at $\varkappa\varepsilon$ and cancelling
the finite tail gives $\int_\varepsilon^{\varkappa\varepsilon} k(u)\,du/u = 0$; on that interval
$k$ is bounded below by $k(\varkappa\varepsilon)$, which must therefore vanish; and since
$\varepsilon$ was arbitrary, reading the conclusion at $\varepsilon = x/\varkappa$ gives
$k(x) = 0$ for every $x > 0$. So $F \equiv 0$ here too, which is excluded. No covering argument
and no passage from almost everywhere to everywhere is needed.

In the canonical gauge $g_{s,t} = F(t\,\cdot) - F(s\,\cdot)$, so $G(\cdot,\omega)$ is strictly
increasing for every $\omega \ne 0$; and if some increment $\mu_{s,t}$ were carried by a lattice
$p\ZZ$ then its cosine transform at $2\pi/p$ would be the integral of the constant $1$, hence
$1$, so $g_{s,t}(2\pi/p) = 0$ at a nonzero frequency, which the clause just proved forbids.
\end{proof}

% draft: Proposition 7.5' clause (3), together with the unimodality claim in the paragraph
% following it. SPLIT OFF from prop:strict-positivity so that the two elementary clauses stay
% out of the trust base.
% twin: hcs prop:volterra-density carries the same Sato absolute-continuity interface
% (ledger A10 there).
\begin{proposition}[regularity of the kernels from the origin]\label{prop:kernel-regularity}
\uses{thm:main-characterization,prop:canonical-gauge,lem:additivity,def:self-decomposable,prop:strict-positivity,lem:selfdecomposable-exponents}
% CHANGED (review 2026-09-09, R27): the \lean tag names two declarations. The interface proper
% is kernel_regularity_law -- A8 and A9 as they are stated in the sources, about a nondegenerate
% symmetric self-decomposable law -- and kernel_regularity is the ASSEMBLY that checks those
% hypotheses for mu_{0,t}, which is [T] work (main_analysis, strict_positivity_strict). Keeping
% them apart is what lets #print axioms of the interface show the ledger entries and nothing
% else. Statement unchanged.
%
% CHANGED (proof of record 2026-09-09, wave 3): the assignment paragraph's account of the
% hypothesis check is rewritten to the machine-checked route
% (SpatialLine.isSelfDecomposable_kernel, SpatialLine.kernel_ne_dirac_zero). Neither of the
% two nodes the R27 note named is on the path. Self-decomposability of mu_{0,t} does not go
% through lem:selfdecomposable-exponents or thm:main-characterization at all: because the gauge
% is ONTO [0,infty), the factor law that def:self-decomposable asks for at ratio b is an
% increment of the family itself, mu_{t',t} with chi(t') = chi(t)/b, and lem:additivity
% supplies its transform. And nondegeneracy needs only that SOME frequency is moved, which is
% prop:canonical-gauge's last clause, not prop:strict-positivity(2). Both old labels are kept
% in \uses, over-approximating being harmless, and prop:canonical-gauge, lem:additivity and
% def:self-decomposable are added. Statement unchanged.
%
% The interface proper is now an admitted Lean axiom, SpatialLine.kernel_regularity_law, on
% blueprint/trust-boundary.txt with A8's and A9's **Lean:** lines naming it; the assembly
% SpatialLine.kernel_regularity is machine-checked and its #print axioms is Lean core plus that
% one name. The node is \leanok in that sense and in no other: what is verified is the
% ASSEMBLY, and the analysis itself remains cited.
%
% CHANGED (fidelity review 2026-09-10, F6-1): the two Lean declarations are NARROWED at the
% origin. Both stated the monotonicity conjunct as AntitoneOn p (Set.Ici 0), on the CLOSED
% half-line; that forces p x <= p 0 < infinity for every x >= 0, i.e. a density bounded by a
% real number, and it is FALSE. Sato Thm. 27.13 gives absolute continuity and Yamazato Thm. 1
% gives convexity left of the mode and concavity right of it; neither pins a value AT the mode,
% and a concave distribution function on (0,infty) may have F'(0+) = +infinity. The
% counterexample is inside this development: the Matern member at gamma <= 1/2 -- admitted by
% matern_exponent and built for every gamma > 0 by witness_two_members_matern_family -- is
% symmetric, nondegenerate and self-decomposable with density of order |x|^(2gamma-1) at the
% origin, unbounded for gamma < 1/2 and logarithmically unbounded at gamma = 1/2 (which is what
% ledger A10 says in Sato's vocabulary, c = 2gamma < 1). Both declarations now read
% AntitoneOn p (Set.Ioi 0), the development's own convention for a profile
% (SDProfile.k_antitone). THE STATEMENT BELOW IS UNCHANGED and remains true: "unimodal with
% mode at the origin" is a claim about the shape of the density and not about its value there,
% and this is exactly the case where prose may leave the reading open and a formal statement may
% not. Recorded at A8 and A9. The other two consumers of the axiom
% (SpatialLine.polya_frequency, SpatialLine.exists_smoothing_law) destructure the
% absolute-continuity conjunct only, so nothing downstream rested on the false clause.
\lean{SpatialLine.kernel_regularity_law, SpatialLine.kernel_regularity}\leanok
Let $(\Phis{s}{t})$ satisfy (A1)--(A8) and (ND). Then for every $t > 0$ the kernel $\mu_{0,t}$
is absolutely continuous, and its density is unimodal with mode at the origin.
\statusA\quad\ledger{A8}\quad\ledger{A9}\quad\emph{(\textbf{Assignment.} \ledger{A8} carries
absolute continuity and \ledger{A9} unimodality, both for nondegenerate self-decomposable laws
on $\RR$. What they do not carry is the hypothesis check, and that check is cheaper than it
looks: that $\mu_{0,t}$ is self-decomposable is \ref{prop:canonical-gauge} with
\ref{lem:additivity} --- the factor law \ref{def:self-decomposable} asks for at ratio $b$ is
the increment $\mu_{t',t}$ at the scale $t'$ with $\chi(t') = \chi(t)/b$, which exists because
the gauge is onto $[0,\infty)$ --- and that it is nondegenerate is the last clause of
\ref{prop:canonical-gauge}, that some frequency is moved. Both are \textbf{[T]} and neither
needs \ref{lem:selfdecomposable-exponents} or \ref{prop:strict-positivity}(2). The symmetric
case places the mode at the origin, which is elementary given unimodality. The statement is
made for the kernels \emph{from the origin} only: the increments $\mu_{s,t}$ with $s > 0$ are
infinitely divisible but not in general self-decomposable, and nothing is claimed about them.)}
\end{proposition}

% ADDED (article mirror 2026-09-14, R157): an [A] node with no proof environment printed nothing
% between its two cited facts and its conclusion, and the hypothesis check --- which is this
% article's work, not the sources' --- sat in the annotation alone. The paper prints the proof
% below; the blueprint is the text of record, so it carries it too. The mode argument is the
% external review's repair: an even unimodal density whose mode is at m /= 0 does NOT have a
% second mode at -m to contradict, because the set of modes is an interval; what the symmetry
% gives is that the interval is symmetric about 0, hence contains it.
\begin{proof}
\leanok
For a nondegenerate self-decomposable law on $\RR$ the statement is the two cited facts:
absolute continuity is \sscite{sato1999levy}, Thm.~27.13, and unimodality is
\sscite{yamazato1978unimodality}, Thm.~1. What remains is the hypothesis check. That
$\mu_{0,t}$ is self-decomposable is \ref{prop:canonical-gauge} with \ref{lem:additivity}: the
factor law \ref{def:self-decomposable} asks for at ratio $b$ is the increment $\mu_{t',t}$ at
the scale $t'$ with $\chi(t') = \chi(t)/b$, which exists because the gauge is onto
$[0,\infty)$. That it is nondegenerate is the last clause of \ref{prop:canonical-gauge}, that
some frequency is moved. The mode is at the origin because the density is even: the set of
modes of a unimodal law is an interval, and for an even density that interval is symmetric
about $0$, so it contains $0$.
\end{proof}

So the Fourier theory ends where the Laplace theory began. The causal vanishing lemma is false
for $\NDs$ (\ref{lem:lattice-zero}) but true, monotonicity in $\omega$ included, for the
exponents the axioms select (\ref{prop:strict-positivity}); and the axioms never produce a
multimodal kernel from the origin. No truncation of the Gaussian is admissible, for a
different reason: the transform of an even density carried by $[-R,R]$ with $p(R-) > 0$ has
the leading term $2p(R-)\sin(R\omega)/\omega$ as $\omega \to \infty$, by integration by parts,
so it changes sign, which \ref{lem:nonvanishing} forbids.

\begin{remark}[the scale evolution, in one line]\label{rem:scale-evolution-preview}
In the canonical gauge $\hat\phi_t(\omega) = e^{-F(t\omega)}$, so on the Fourier side the
family solves $\partial_t\hat u = -t^{-1}B(t\omega)\,\hat u$ with $B = \omega F'$ the symbol of
\ref{lem:selfdecomposable-exponents}(2). By \ref{eq:symbol} that symbol is the generator of the
symmetric \Levy{} process with Gaussian part $2a$ and jump measure $\varpi$. For the Gaussian,
$B(\omega) = 2a\omega^2$ and the equation is the heat equation in the variance gauge; for the
\Matern{} member, $F = \gamma\log(1+\omega^2)$ and $B(\omega) = 2\gamma\omega^2/(1+\omega^2)$, a
\emph{bounded} symbol. That the
symbol be a polynomial, so that the evolution is a differential equation in $x$, forces
$\varpi = 0$ by \ref{eq:symbol}. The equation and the observation are stated here without
proof; the evolution in generator form is Chapter~\ref{ch:generator} and the locality theorem
they lead to is Chapter~\ref{ch:locality}.
\end{remark}
